\chapter*{Abstract}

This report presents the design, implementation, and evaluation of a dual-mode multimodal book recommendation system built on a data platform with modular ingestion, storage, and retrieval components, backed by MongoDB for persistent interaction logs and Redis as an asynchronous interaction write-queue.
The system addresses two distinct recommendation scenarios: proactive candidate generation based on long-term behavioural history, and intent-sensitive ranking based on recent interaction sequences.

Two complementary pipelines are proposed.
Pipeline~A retrieves candidates through a Cleora co-purchase graph embedding index covering 375{,}280 items, then re-ranks them using a BGE-M3 semantic profile vector constructed by mean-pooling the user's interaction embeddings.
Pipeline~B applies a DIF-SASRec transformer with decoupled content and category attention streams, pretrained on 100{,}000 users, to model sequential reading intent.
A content veto filter (cosine threshold $\tau = 0.3$) removes semantically unrelated candidates before re-ranking.
The final recommendations are formed by the union of the top-10 outputs of both pipelines.

The system is evaluated on the Amazon Book Reviews dataset using a leave-one-out protocol with sampled negatives.
On 100{,}000 test users with $N = 99$ negatives, the combined system achieves \textbf{HR@10 = 0.9736} and \textbf{NDCG@10 = 0.5571} — a 124\% improvement in hit rate over the content-only baseline and an improvement of more than $8\times$ over the prior GRU-SeqDQN approach.
Evaluation across five random seeds confirms stability ($\sigma = 0.0024$ for HR@10).

An active text search pipeline supports semantic retrieval over 3 million items via FAISS HNSW approximate search, BM25 re-ranking through RRF fusion, and NLLB-based Vietnamese query translation.
Latency optimisations reduce end-to-end query response time from 1{,}102 ms to 59 ms (warm cache), an 18.6$\times$ improvement.
BGE-M3 achieves up to 11$\times$ higher Precision@10 than the legacy BLaIR encoder on Vietnamese queries, validating the choice of a multilingual encoder for the target user population.


